\documentclass[tikz,border=6pt]{standalone}
\usepackage{ctex}
\usepackage{amsmath}
\usetikzlibrary{calc}

\begin{document}
\begin{tikzpicture}[
  x={(1cm,0cm)}, y={(0,1cm)}, z={(-0.5cm,-0.4cm)},
  line cap=round, line join=round,
  every label/.style={font=\small}
]

% 正方体 8 种截面 — 4 种典型形态（2×2 排列）
% 棱长 = 2，坐标范围 [0,2]^3

\newcommand{\drawcubegray}{
  \coordinate (A) at (0,0,0);
  \coordinate (B) at (2,0,0);
  \coordinate (C) at (2,0,2);
  \coordinate (D) at (0,0,2);
  \coordinate (Ap) at (0,2,0);
  \coordinate (Bp) at (2,2,0);
  \coordinate (Cp) at (2,2,2);
  \coordinate (Dp) at (0,2,2);
  % 后隐棱（虚线）
  \draw[gray!50, dashed, thin] (A) -- (Ap);
  \draw[gray!50, dashed, thin] (A) -- (B);
  \draw[gray!50, dashed, thin] (A) -- (D);
  % 可见棱
  \draw[gray!70, thick] (B) -- (C) -- (D);
  \draw[gray!70, thick] (Ap) -- (Bp) -- (Cp) -- (Dp) -- cycle;
  \draw[gray!70, thick] (B) -- (Bp);
  \draw[gray!70, thick] (C) -- (Cp);
  \draw[gray!70, thick] (D) -- (Dp);
}

%========== 行1 左：三角形截面（过顶点 A(0,0,0), B(2,0,0), Dp(0,2,2)）==========
\begin{scope}[xshift=0cm, yshift=0cm]
  \drawcubegray
  \coordinate (A) at (0,0,0);
  \coordinate (B) at (2,0,0);
  \coordinate (Dp) at (0,2,2);
  % 截面（红色三角形）
  \fill[red!25, opacity=0.7] (A) -- (B) -- (Dp) -- cycle;
  \draw[red, very thick] (A) -- (B) -- (Dp) -- cycle;
  % 顶点标注
  \node[red, below left, font=\scriptsize] at (A) {$A$};
  \node[red, below right, font=\scriptsize] at (B) {$B$};
  \node[red, above left, font=\scriptsize] at (Dp) {$D'$};
  % 小标题
  \node[font=\small\bfseries, red] at (1,-1.5,1) {(1) 三角形截面};
  \node[font=\scriptsize, gray] at (1,-1.95,1) {过 $A,B,D'$ 三顶点};
\end{scope}

%========== 行1 右：矩形截面（平行于正面的中截面）==========
\begin{scope}[xshift=5cm, yshift=0cm]
  \drawcubegray
  % 中截面：y=1 处（平行于 xOz 面）
  \coordinate (M1) at (0,1,0);
  \coordinate (M2) at (2,1,0);
  \coordinate (M3) at (2,1,2);
  \coordinate (M4) at (0,1,2);
  \fill[blue!25, opacity=0.7] (M1) -- (M2) -- (M3) -- (M4) -- cycle;
  \draw[blue, very thick] (M1) -- (M2) -- (M3) -- (M4) -- cycle;
  % 中点标注
  \fill[blue] (M1) circle [radius=1.8pt];
  \fill[blue] (M2) circle [radius=1.8pt];
  \fill[blue] (M3) circle [radius=1.8pt];
  \fill[blue] (M4) circle [radius=1.8pt];
  % 标尺：显示是正方形/矩形
  \node[font=\small\bfseries, blue] at (1,-1.5,1) {(2) 矩形截面};
  \node[font=\scriptsize, gray] at (1,-1.95,1) {平行于某面的中间截面};
\end{scope}

%========== 行2 左：五边形截面（斜切，不对称）==========
\begin{scope}[xshift=0cm, yshift=-5.5cm]
  \drawcubegray
  % 五边形截面（过 5 个点）
  % 截面由平面 x + y + z = 4 与正方体 [0,2]^3 的交截
  % 计算交点：
  % 棱 AB: y=0,z=0 → x=4 (超出，无交)
  % 棱 BC: x=2,z=0 → 2+y=4 → y=2，即点 (2,2,0)=Bp √
  % 棱 CD: x=2,y=0 → 2+z=4 → z=2，即点 (2,0,2)=C √
  % 棱 DDp: x=0,y=0 → z=4 (超出)
  % 棱 BpCp: x=2,y=2 → z=0，即 Bp=(2,2,0) 已有
  % 棱 CpDp: y=2,z=2 → x=0，即 (0,2,2)=Dp √
  % 棱 DpAp: x=0,z=2 → y=2，即 Dp 已有
  % 棱 ApBp: z=0,x=0 → y=4 超出
  % 棱 ADp 用不同平面：改用 x+0.5y+z=3
  % 棱 AB(y=0,z=0): x=3>2，无
  % 棱 BC(x=2,z=0): 2+0.5y=3→y=2，B'=(2,2,0)=Bp
  % 棱 BBp(x=2,z=0,y∈[0,2]): 2+0.5y=3→y=2，Bp
  % 棱 CD(x=2,y=0,z∈[0,2]): 2+z=3→z=1，E=(2,0,1)
  % 棱 DDp(x=0,y=0,z∈[0,2]): z=3>2，无
  % 棱 DpCp(y=2,z=2): 0.5*2+x+2=3→x=-1<0，无
  % 棱 Dp Ap(x=0,z=2): 0.5y+2=3→y=2，Dp=(0,2,2)  √
  % 棱 BpCp(x=2,y=2): 2+1+z=3→z=0，Bp 已有
  % 棱 ApBp(y=2,z=0): x+1=3→x=2，Bp 已有
  % 棱 Cp Dp(y=2,z=2): x+1+2=3→x=0，Dp 已有
  % 棱 ApA(x=0,y∈[0,2],z=0): 0.5y=3>2，无
  % 棱 ADp 侧：取平面 x+y+1.5z=4
  % 棱 AB(y=0,z=0):x=4>2无
  % 棱 BC(x=2,z=0):y=2→Bp=(2,2,0)
  % 棱 CD(x=2,y=0):2+1.5z=4→z=4/3, F=(2,0,4/3)
  % 棱 DDp(x=0,y=0):1.5z=4→z=8/3>2无
  % 棱 DCp(x=0,z=2):y+3=4→y=1, G=(0,1,2)  [棱DpAp是x=0,z=2]
  % 棱 DpAp(x=0,z=2,y∈[0,2]):y+3=4→y=1, G=(0,1,2)√
  % 棱 ApBp(y=2,z=0):x+2=4→x=2，Bp已有
  % 棱 BpCp(x=2,y=2):2+2+1.5z=4→z=-2/3<0，无
  % 棱 CpDp(y=2,z=2):x+2+3=4→x=-1<0无
  % 棱 AApDpD边→ApAp端：棱AAp(x=0,z=0,y∈[0,2]):y=4>2无
  % 所以截面为 Bp-F-G 三点 → 三角形，不对称5边需调整
  % 换截面：x + y + 1.2z = 3.5
  % 棱BC(x=2,z=0):y=1.5,H=(2,1.5,0)
  % 棱CD(x=2,y=0):2+1.2z=3.5→z=1.25,I=(2,0,1.25)
  % 棱DDp(x=0,y=0):1.2z=3.5→z=2.917>2无
  % 棱DpAp(x=0,z=2,y∈[0,2]):y+2.4=3.5→y=1.1,J=(0,1.1,2)
  % 棱CpDp(y=2,z=2):x+2+2.4=3.5→x=-0.9<0无
  % 棱BpCp(x=2,y=2):2+2+1.2z=3.5→z=-0.42<0无
  % 棱ApBp(y=2,z=0):x+2=3.5→x=1.5,K=(1.5,2,0)
  % 棱AAp(x=0,y∈[0,2],z=0):y=3.5>2无
  % 棱ABp对角?不是正方体的棱
  % 5点: H=(2,1.5,0), I=(2,0,1.25), J=(0,1.1,2), K=(1.5,2,0), 还需第5点
  % 棱DpCp: y=2,z=2: x+4.4=3.5→x<0
  % 棱Dp-Ap: done→J
  % 棱A-D: x=0,y=0,z∈[0,2]: 1.2z=3.5>2无
  % 棱B-C: done→I
  % 五个点: H,I,J,K，还缺一个
  % 棱Cp-Bp: done，无
  % 检查棱 A(0,0,0)-Ap(0,2,0)无；棱D(0,0,2)-C(2,0,2)
  %   z=2,y=0: x+0+2.4=3.5→x=1.1, L=(1.1,0,2)√
  % 5点顺序：H=(2,1.5,0), K=(1.5,2,0), J=(0,1.1,2), L=(1.1,0,2), I=(2,0,1.25)
  \coordinate (PH) at (2,1.5,0);
  \coordinate (PI) at (2,0,1.25);
  \coordinate (PJ) at (0,1.1,2);
  \coordinate (PK) at (1.5,2,0);
  \coordinate (PL) at (1.1,0,2);
  \fill[green!25, opacity=0.7] (PH) -- (PK) -- (PJ) -- (PL) -- (PI) -- cycle;
  \draw[green!60!black, very thick] (PH) -- (PK) -- (PJ) -- (PL) -- (PI) -- cycle;
  \foreach \pt in {PH,PI,PJ,PK,PL}{
    \fill[green!60!black] (\pt) circle [radius=1.8pt];
  }
  \node[font=\small\bfseries, green!60!black] at (1,-1.5,1) {(3) 五边形截面};
  \node[font=\scriptsize, gray] at (1,-1.95,1) {斜切，不过顶点};
\end{scope}

%========== 行2 右：六边形截面（过6棱中点 → 正六边形）==========
\begin{scope}[xshift=5cm, yshift=-5.5cm]
  \drawcubegray
  % 六边形：正方体 6 条棱中点
  % 平面 x+y+z=3（过 (1,0,2),(2,1,0),(0,2,1),(1,2,0),(2,0,1),(0,1,2) 中点）
  % 棱AB中(y=0,z=0): x=3>2 无
  % 用平面 x+y+z=3的正方体[0,2]^3截面：
  % 棱BC(x=2,z=0): y=1, Q1=(2,1,0)
  % 棱CD(x=2,y=0): z=1, Q2=(2,0,1)
  % 棱DDp(x=0,y=0): z=3>2 无
  % 棱DC: done
  % 棱ADp: x=0,y↑,z=2: 棱DpAp(x=0,z=2): y=1, Q3=(0,1,2)
  % 棱DpCp: y=2,z=2: x=-1<0 无
  % 棱Dp-D: z=2,x=0,y↑: 棱DpAp done
  % 棱Cp-Dp: done无
  % 棱Cp-Bp: x=2,y=2: z=-1<0 无
  % 棱Ap-Bp: y=2,z=0: x=1, Q4=(1,2,0)
  % 棱ApAp: 就是Ap=Q4
  % 棱Ap-A: x=0,z=0,y: y=3>2无
  % 棱A-D(x=0,y=0,z∈[0,2]): z=3>2无
  % 棱BpCp(x=2,y=2): z=-1<0无
  % 棱CpDp(y=2,z=2): x=-1<0无
  % 只找到 Q1,Q2,Q3,Q4 → 需要其他棱
  % 棱B(2,0,0)-Bp(2,2,0): x=2,z=0,y∈[0,2]: y=1→Q1已有
  % 棱Bp(2,2,0)-Cp(2,2,2): x=2,y=2,z∈: z=-1无
  % 棱A(0,0,0)-Ap(0,2,0): x=0,z=0,y: y=3无
  % 棱D(0,0,2)-C(2,0,2): y=0,z=2,x∈: x=1, Q5=(1,0,2)
  % 棱Dp(0,2,2)-Cp(2,2,2): y=2,z=2,x∈: x=-1无
  % 棱C(2,0,2)-Cp(2,2,2): x=2,z=2,y∈: y=-1无
  % 棱Ap(0,2,0)-Dp(0,2,2): x=0,y=2,z∈: z=1, Q6=(0,2,1)
  % 六点: Q1=(2,1,0), Q2=(2,0,1), Q5=(1,0,2), Q3=(0,1,2), Q6=(0,2,1), Q4=(1,2,0)
  \coordinate (Q1) at (2,1,0);
  \coordinate (Q2) at (2,0,1);
  \coordinate (Q3) at (0,1,2);
  \coordinate (Q4) at (1,2,0);
  \coordinate (Q5) at (1,0,2);
  \coordinate (Q6) at (0,2,1);
  \fill[purple!25, opacity=0.7] (Q1) -- (Q2) -- (Q5) -- (Q3) -- (Q6) -- (Q4) -- cycle;
  \draw[purple, very thick] (Q1) -- (Q2) -- (Q5) -- (Q3) -- (Q6) -- (Q4) -- cycle;
  \foreach \pt in {Q1,Q2,Q3,Q4,Q5,Q6}{
    \fill[purple] (\pt) circle [radius=1.8pt];
  }
  \node[font=\small\bfseries, purple] at (1,-1.5,1) {(4) 六边形截面（正六边形）};
  \node[font=\scriptsize, gray] at (1,-1.95,1) {过各棱中点，$x{+}y{+}z{=}3$};
\end{scope}

% 总标题
\node[font=\large\bfseries] at (4.5, 2.8) {正方体的典型截面};

\end{tikzpicture}
\end{document}
